\documentclass[10pt,a4paper]{amsart}
\usepackage{amsmath,amssymb,amsthm,graphicx}
\usepackage{verbatim}
\newenvironment{reference}{\comment}{\endcomment}
%\newenvironment{reference}{}{}
\newenvironment{slogan}{\comment}{\endcomment}
\newenvironment{history}{\comment}{\endcomment}

% For commutative diagrams we use Xy-pic
\usepackage[all]{xy}

% We use 2cell for 2-commutative diagrams.
\xyoption{2cell}
% Theorem environments.
%
\theoremstyle{plain}
\newtheorem{theorem}[subsection]{Theorem}
\newtheorem{proposition}[subsection]{Proposition}
\newtheorem{lemma}[subsection]{Lemma}

\theoremstyle{definition}
\newtheorem{definition}[subsection]{Definition}
\newtheorem{example}[subsection]{Example}
\newtheorem{exercise}[subsection]{Exercise}
\newtheorem{situation}[subsection]{Situation}

\theoremstyle{remark}
\newtheorem{remark}[subsection]{Remark}
\newtheorem{remarks}[subsection]{Remarks}

\numberwithin{equation}{subsection}

% Macros
%
\def\lim{\mathop{\mathrm{lim}}\nolimits}
\def\colim{\mathop{\mathrm{colim}}\nolimits}
\def\Spec{\mathop{\mathrm{Spec}}}
\def\Hom{\mathop{\mathrm{Hom}}\nolimits}
\def\Ext{\mathop{\mathrm{Ext}}\nolimits}
\def\SheafHom{\mathop{\mathcal{H}\!\mathit{om}}\nolimits}
\def\SheafExt{\mathop{\mathcal{E}\!\mathit{xt}}\nolimits}
\def\Sch{\mathit{Sch}}
\def\Mor{\mathop{\mathrm{Mor}}\nolimits}
\def\Ob{\mathop{\mathrm{Ob}}\nolimits}
\def\Sh{\mathop{\mathit{Sh}}\nolimits}
\def\NL{\mathop{N\!L}\nolimits}
\def\CH{\mathop{\mathrm{CH}}\nolimits}
\def\proetale{{pro\text{-}\acute{e}tale}}
\def\etale{{\acute{e}tale}}
\def\QCoh{\mathit{QCoh}}
\def\Ker{\mathop{\mathrm{Ker}}}
\def\Im{\mathop{\mathrm{Im}}}
\def\Coker{\mathop{\mathrm{Coker}}}
\def\Coim{\mathop{\mathrm{Coim}}}

% Boxtimes
%
\DeclareMathSymbol{\boxtimes}{\mathbin}{AMSa}{"02}

%
% Macros for moduli stacks/spaces
%
\def\QCohstack{\mathcal{QC}\!\mathit{oh}}
\def\Cohstack{\mathcal{C}\!\mathit{oh}}
\def\Spacesstack{\mathcal{S}\!\mathit{paces}}
\def\Quotfunctor{\mathrm{Quot}}
\def\Hilbfunctor{\mathrm{Hilb}}
\def\Curvesstack{\mathcal{C}\!\mathit{urves}}
\def\Polarizedstack{\mathcal{P}\!\mathit{olarized}}
\def\Complexesstack{\mathcal{C}\!\mathit{omplexes}}
% \Pic is the operator that assigns to X its picard group, usage \Pic(X)
% \Picardstack_{X/B} denotes the Picard stack of X over B
% \Picardfunctor_{X/B} denotes the Picard functor of X over B
\def\Pic{\mathop{\mathrm{Pic}}\nolimits}
\def\Picardstack{\mathcal{P}\!\mathit{ic}}
\def\Picardfunctor{\mathrm{Pic}}
\def\Deformationcategory{\mathcal{D}\!\mathit{ef}}
\usepackage{cleveref}
\crefname{originalreference}{original reference}{original references}
\Crefname{originalreference}{Original reference}{Original references}
\makeatletter
\newlabel{example-glue-points}{{0C1J}{source}{}{sourcepointer-1}{}}
\newlabel{example-glue-points@cref}{{[originalreference][0][]0C1J}{[1][0][]source}{}{}{}}
\makeatother
\crefname{originalreference}{original reference}{original references}
\Crefname{originalreference}{Original reference}{Original references}
\makeatletter
\newlabel{example-squish-tangent-vector}{{0C1K}{source}{}{sourcepointer-2}{}}
\newlabel{example-squish-tangent-vector@cref}{{[originalreference][0][]0C1K}{[1][0][]source}{}{}{}}
\makeatother
\crefname{originalreference}{original reference}{original references}
\Crefname{originalreference}{Original reference}{Original references}
\makeatletter
\newlabel{lemma-no-in-between-over-k}{{0C1I}{source}{}{sourcepointer-3}{}}
\newlabel{lemma-no-in-between-over-k@cref}{{[originalreference][0][]0C1I}{[1][0][]source}{}{}{}}
\makeatother
\title{A finite morphism that is an isomorphism off finitely many points factors into glueings of pairs of points and squishings of tangent vectors}
\author{The Stacks Project Authors}
\date{stacks-project:curves.tex:local-snapshot-20261012; GNU FDL 1.2 or later}
\begin{document}
\maketitle
\noindent\textit{Editorial note.} All delivered mathematical text is the author's own, copied from the pinned Stacks Project chapter file curves.tex; no statement or proof step was rewritten, reordered or removed.\par\medskip
\noindent\textit{Editorial note.} The counted proof is the author's complete argument: the skyscraper decomposition, the induction framework with intermediate subalgebras, the multi-point splitting, the Nakayama reduction at a single point, and the terminal k x k versus k[epsilon] classification are all retained verbatim.\par\medskip
\noindent\textit{Editorial note.} The label lemma-factor-almost-isomorphism names the counted statement, and the proof environment immediately following it in the source is delivered verbatim as its proof.\par\medskip
\noindent\textit{Editorial note.} The counted statement is carried by the author's own lemma environment in the chapter source.\par\medskip
\noindent\textit{Allowed standard prerequisites (editorial).}\begin{itemize}
\item The two elementary modifications. Same chapter, Examples example-glue-points and example-squish-tangent-vector. Define the glueing of two points and the squishing of a tangent vector, the two step types of the factorization.
\item Rings with exactly two subalgebras. Same chapter, Lemma lemma-no-in-between-over-k. Classifies a finite k-algebra with exactly two k-subalgebras as k x k or k[epsilon], the two terminal cases of the induction.
\end{itemize}\medskip
\ifdefined\hypertarget\hypertarget{sourcepointer-1}{}\fi
\noindent\textit{Source reference.} Original reference 0C1J: sources/stacks-curves/pinned-original-source.tex, source lines 3202--3202. Prerequisite pointer invoked at proof line(s) 3267; it names a standard fact of the same chapter that the counted argument uses, and every step of the counted proof built on it is retained verbatim in the delivered body.\par\medskip
\ifdefined\hypertarget\hypertarget{sourcepointer-2}{}\fi
\noindent\textit{Source reference.} Original reference 0C1K: sources/stacks-curves/pinned-original-source.tex, source lines 3228--3228. Prerequisite pointer invoked at proof line(s) 3268; it names a standard fact of the same chapter that the counted argument uses, and every step of the counted proof built on it is retained verbatim in the delivered body.\par\medskip
\ifdefined\hypertarget\hypertarget{sourcepointer-3}{}\fi
\noindent\textit{Source reference.} Original reference 0C1I: sources/stacks-curves/pinned-original-source.tex, source lines 3183--3183. Prerequisite pointer invoked at proof line(s) 3317; it names a standard fact of the same chapter that the counted argument uses, and every step of the counted proof built on it is retained verbatim in the delivered body.\par\medskip
\begin{lemma}
\label{lemma-factor-almost-isomorphism}
Let $k$ be an algebraically closed field. Let $f : X' \to X$ be a
finite morphism of algebraic $k$-schemes such that
$\mathcal{O}_X \subset f_*\mathcal{O}_{X'}$ and such that $f$ is an
isomorphism away from a finite set of points. Then there is a factorization
$$
X' = X_n \to X_{n - 1} \to \ldots \to X_1 \to X_0 = X
$$
such that each $X_i \to X_{i - 1}$ is either the glueing of
two points or the squishing of a tangent vector
(see Examples \ref{example-glue-points} and
\ref{example-squish-tangent-vector}).
\end{lemma}

\begin{proof}
Let $U \subset X$ be the maximal open set over which $f$ is an isomorphism.
Then $X \setminus U = \{x_1, \ldots, x_n\}$ with $x_i \in X(k)$.
We will consider factorizations $X' \to Y \to X$ of $f$ such that
both morphisms are finite and
$$
\mathcal{O}_X \subset g_*\mathcal{O}_Y \subset f_*\mathcal{O}_{X'}
$$
where $g : Y \to X$ is the given morphism. By assumption
$\mathcal{O}_{X, x} \to (f_*\mathcal{O}_{X'})_x$ is an isomorphism
onless $x = x_i$ for some $i$. Hence the cokernel
$$
f_*\mathcal{O}_{X'}/\mathcal{O}_X = \bigoplus \mathcal{Q}_i
$$
is a direct sum of skyscraper sheaves $\mathcal{Q}_i$ supported at
$x_1, \ldots, x_n$.
Because the displayed quotient is a coherent $\mathcal{O}_X$-module,
we conclude that $\mathcal{Q}_i$ has finite length over
$\mathcal{O}_{X, x_i}$. Hence we can argue
by induction on the sum of these lengths, i.e., the length of
the whole cokernel.

\medskip\noindent
If $n > 1$, then we can define an $\mathcal{O}_X$-subalgebra
$\mathcal{A} \subset f_*\mathcal{O}_{X'}$ by taking the inverse
image of $\mathcal{Q}_1$. This will give a nontrivial factorization
and we win by induction.

\medskip\noindent
Assume $n = 1$. We abbreviate $x = x_1$. Consider the finite
$k$-algebra extension
$$
A = \mathcal{O}_{X, x} \subset (f_*\mathcal{O}_{X'})_x = B
$$
Note that $\mathcal{Q} = \mathcal{Q}_1$ is the skyscraper sheaf
with value $B/A$.
We have a $k$-subalgebra $A \subset A + \mathfrak m_A B \subset B$.
If both inclusions are strict, then we obtain a nontrivial
factorization and we win by induction as above.
If $A + \mathfrak m_A B = B$, then $A = B$ by Nakayama, then
$f$ is an isomorphism and there is nothing to prove.
We conclude that we may assume $B = A + \mathfrak m_A B$.
Set $C = B/\mathfrak m_A B$. If $C$ has more than $2$
$k$-subalgebras, then we obtain a subalgebra between $A$
and $B$ by taking the inverse image in $B$. Thus we may
assume $C$ has exactly $2$ $k$-subalgebras. Thus $C = k \times k$
or $C = k[\epsilon]$ by Lemma \ref{lemma-no-in-between-over-k}.
In this case $f$ is correspondingly the glueing two points or the
squishing of a tangent vector.
\end{proof}

\end{document}
